\documentclass[aspectratio=169,11pt]{beamer}

\usepackage[T1]{fontenc}
\usepackage[utf8]{inputenc}
\usepackage{lmodern}
\usepackage{amsmath,amssymb}
\usepackage{booktabs}
\usepackage{xcolor}
\usepackage{hyperref}
\usepackage{microtype}

\usetheme{Madrid}
\usecolortheme{default}

\definecolor{LaurierPurple}{HTML}{4B116F}
\definecolor{DeepBlue}{HTML}{0047AB}

\setbeamertemplate{navigation symbols}{}
\setbeamercolor{title}{fg=white,bg=LaurierPurple}
\setbeamercolor{frametitle}{fg=white,bg=LaurierPurple}
\setbeamercolor{block title}{fg=white,bg=DeepBlue}
\setbeamercolor{structure}{fg=DeepBlue}

\hypersetup{colorlinks=true,urlcolor=DeepBlue,linkcolor=DeepBlue}

\title[Applied Project]{Module 9: Applied Machine Learning and Causal Project}
\subtitle{EC 410K / EC 610I --- Machine Learning in Economics}
\author{M. Jahangir Alam, PhD}
\institute{Wilfrid Laurier University}
\date{Spring 2026}

\begin{document}

\begin{frame}[plain]
  \titlepage
\end{frame}

\begin{frame}{Module overview}
  \begin{itemize}
    \item The group project asks: \textbf{given a replicated empirical paper, what can ML add}?
    \item Successful projects are explicit about the \textbf{task} (prediction, heterogeneity, measurement, targeting) and the \textbf{evaluation}.
    \item The best work connects ML outputs to economic interpretation without confusing correlation with causation.
  \end{itemize}
\end{frame}

\begin{frame}{Learning objectives}
  \begin{enumerate}
    \item Draft a feasible ML extension that matches the paper’s data and question.
    \item Choose baselines and tune models with honest validation.
    \item Communicate limitations: external validity, sample size, endogeneity, overfitting.
    \item Produce outputs aligned with the final report + poster.
  \end{enumerate}
\end{frame}

\begin{frame}{Three common project templates}
  \begin{enumerate}
    \item \textbf{Prediction upgrade:} can ML improve out-of-sample prediction of a key outcome?
    \item \textbf{Heterogeneity:} do effects vary across units in economically interpretable ways?
    \item \textbf{Measurement / text:} can new variables be extracted to enrich the analysis?
  \end{enumerate}
\end{frame}

\begin{frame}{What reviewers will look for}
  \begin{itemize}
    \item Clear link between replication baseline and ML extension.
    \item Correct metrics for the stated goal (not only in-sample \(R^2\)).
    \item Interpretability where appropriate (coefficients, SHAP, partial dependence---used carefully).
    \item Reproducible code and data documentation.
  \end{itemize}
\end{frame}

\begin{frame}{Avoid these project failures}
  \begin{itemize}
    \item ``We ran random forest'' without a stated economic question.
    \item Training on the full sample and reporting training accuracy as success.
    \item Claiming causality from predictive feature importance.
    \item ML extension unrelated to the paper’s setting or outcomes.
  \end{itemize}
\end{frame}

\begin{frame}{Python / Colab demonstration}
  \begin{itemize}
    \item Dry-run comparison: Ridge baseline vs random forest on a nonlinear simulated surface.
    \item Use the pattern as a coding template: split, tune, evaluate, summarize.
  \end{itemize}
\end{frame}

\begin{frame}[fragile]{Colab setup}
  \begin{verbatim}
!pip -q install numpy pandas scikit-learn
  \end{verbatim}
\end{frame}

\begin{frame}{Student / group assignment}
  \begin{itemize}
    \item Peer review another team’s 1-page plan using a rubric: question, baseline, metric, risks.
    \item Revise plan; run the dry-run notebook; begin implementing on replication data.
    \item Milestone: minimal viable result (one clean table or figure).
  \end{itemize}
\end{frame}

\begin{frame}{Summary}
  \begin{itemize}
    \item Integration beats complexity: one sharp extension beats many shallow models.
    \item Tie every result to economics and to what the original paper already established.
    \item Build the poster storyline early: motivation \(\rightarrow\) replication \(\rightarrow\) ML \(\rightarrow\) takeaway.
  \end{itemize}
\end{frame}

\begin{frame}[plain]
  \begin{center}
    \Large Questions?\\[0.8em]
    \normalsize Next: research communication (Module 10).
  \end{center}
\end{frame}

\end{document}
